\chapter{La Caldera como infraestructura}

\begin{cautela}
Lo que este capítulo lee en el presente, el capítulo \ref{cap:presencia} lo encontró ya escrito entre 1923 y 1945: policía, registro, recaudación y vías de paso, y casi nada más. Que el reparto sea el mismo ochenta años después es el argumento de fondo de estas páginas, y conviene decir que no lo produjo esta lectura: lo produjo el barrido edición por edición de aquel tramo.
\end{cautela}

\label{cap:infraestructura}
\begin{cautela}
Este capítulo es distinto de todos los anteriores y conviene decirlo antes de empezar. Sobre todo relee lo documentado. Agrega pocas piezas nuevas, que se identifican como tales: las declaraciones vecinales de enero de 2026, los boletines municipales de la ciudad de Salta y un grupo de actos de 1935 a 1945 ---los Decretos 4217/40, 3016-H/44, 6749-H/42 y 5738-G/45, el informe sobre las truchas y los accidentes de Gallinato--- que sólo se desarrollan aquí. Todo lo que afirma se apoya en material que los
capítulos previos establecieron con fuente, pero la afirmación misma —lo que ese material
significa junto— es una lectura del autor y no un hallazgo.

Va hacia el final porque es lo que queda cuando uno cierra el libro y se pregunta qué le pasa a
este lugar. Y va marcado, como todo lo inferido en este trabajo, para que se lo pueda
discutir sin discutir los documentos.
\end{cautela}

\section{Las tesis describen cómo; falta el qué}

Las tres tesis de este libro describen mecanismos. La primera dice que los dos circuitos de
producción de suelo están regulados en proporción inversa a la capacidad de sus
destinatarios. La segunda, que la opacidad no está en el acto sino en el criterio. La
tercera, que el dispositivo es anterior al mercado que lo usa.

Las tres son verificables y ninguna requiere que nadie haya decidido producir el resultado.
Pero las tres contestan \textit{cómo} funciona esto. Ninguna contesta la pregunta que un
lector de La Caldera haría primero, que es más simple: \textit{qué le pasa a este lugar}.

La respuesta que se desprende del material, tomado en conjunto, cabe en una frase.

\begin{quote}
\textbf{La Caldera es infraestructura de una ciudad que no la gobierna.}
\end{quote}

\section{Qué sale del departamento}

Póngase en una sola página lo que los capítulos documentaron por separado.

\textbf{Sale el agua.} En junio de 1972 un gobierno de facto expropió seiscientas ochenta y
seis hectáreas para construir el embalse Campo Alegre, con toma de posesión previa a la
discusión del precio. El acueducto que lleva esa agua a la cisterna de El Huaico recorre
unos veinte kilómetros y atraviesa los dos municipios del departamento. La microcuenca cuyo
bosque se remueve para lotear es la zona de recarga del sistema que abastece a la ciudad de
Salta.

\textbf{Sale la piedra.} Entre 2014 y 2015 este libro documentó, por boletines sueltos, seis
concesiones de áridos sobre la cuenca; tres de ellas sobre tierra fiscal, y dos de esas tres ---más de treinta hectáreas--- en manos de una misma firma. La nómina oficial que la autoridad minera elevó al
juzgado en 2022 --- capítulo~\ref{cap:resistencias} --- lleva ese número a
\textbf{veinte concesiones vigentes}, de las cuales cinco tienen la evaluación ambiental
aprobada. El material no se usa en el departamento: alimenta la obra de la capital y de la
región.

\textbf{Sale el suelo, aunque no se mueva.} La casa quinta de Lesser rematada en 1927
salió a subasta junto con derechos sobre una casa en la capital; de la de 1929, rematada en
Vaqueros, el aviso no dice dónde vivía su dueño. Casi un siglo después,
una de cada cinco viviendas del departamento no tiene residentes permanentes —el doble de
la proporción provincial, y en una de sus tres fracciones censales, más del cuarenta por
ciento—. Y esa tasa ya era del diecinueve por ciento en 2010, antes del ciclo de loteos que
este libro documenta.

\textbf{Sale el fin de semana.} En Carnaval de 2026 el municipio informó una ocupación
turística del ochenta por ciento, con cuatro establecimientos completos. Dos de los hoteles
que hicieron posible esa capacidad recibieron, en 2014 y 2015, exención de tributos
provinciales; y uno de ellos, además, cuarenta y una hectáreas de tierra fiscal en comodato
gratuito por veinte años para una cancha de golf.

\textbf{Y sale hasta la planificación.} El plan urbano de Vaqueros lo redactó una consultora seleccionada por una unidad ejecutora provincial, con financiamiento del Banco Interamericano de Desarrollo y no objeción de una unidad ejecutora nacional; el de La Caldera lo encargó el Consejo Federal de Inversiones a la Fundación del Colegio de Arquitectos y a la Subsecretaría provincial de Planeamiento. En los dos casos el municipio recibió el plan; no lo encargó ni lo pagó. El informe final del de
La Caldera nombra, en su primer párrafo, a un municipio situado a cuatrocientos kilómetros.

\section{Quién decide}

Ahora póngase al lado la otra columna.

En 1908 el municipio no recaudaba: \textbf{vendía en licitación el derecho a cobrar sus
impuestos}, y casi el ochenta por ciento de su presupuesto anual salía de esa venta. Un municipio
que remata su recaudación no necesita saber cuántos inmuebles tiene ni cuánto valen.

\textbf{Y el Estado provincial gastaba en un solo cargo del departamento casi lo que el
municipio entero.} En 1912 el subcomisario interino del pueblo cobraba cien pesos mensuales:
\$1.200 al año, cuatro quintos de los \$1.483,77 de la renta municipal de 1908. En 1934 la
escuela provincial del pueblo costaba \$3.900 al año, más de una vez y un tercio los dos mil
ochocientos y tantos que el capítulo \ref{cap:hacienda} infiere como renta municipal de 1937
(capítulos \ref{cap:presencia} y \ref{cap:hacienda}). \textbf{Un municipio cuya renta anual
equivale al sueldo de un subcomisario no tiene con qué regular el suelo de nadie}, que es lo que
la advertencia preliminar llama asignarle esa función «sin darle con qué».

En 1922 el Interventor Nacional \textbf{designó a los cinco miembros} de su comisión
municipal, y tres meses más tarde el Poder Ejecutivo provincial \textbf{le aprobó el
presupuesto} previo dictamen del Fiscal General.

En 1983 seguía siendo una Comisión Municipal cuyo presidente concentraba las funciones del
cuerpo deliberativo por la Ley de facto 5686, y cuyas ordenanzas —incluida una rebaja de
tasas para jubilados— necesitaban aprobación de la Secretaría de Estado de Municipalidades.

En 2022, con el departamento en doce mil doscientos noventa y nueve habitantes tras crecer un cincuenta y ocho por ciento, la estadística provincial sigue clasificando a sus dos municipios en \textbf{tercera categoría} ---la de las comisiones municipales de una ley orgánica
que se derogó en 2018 y cuyo reemplazo ya no tiene categorías---.

\textbf{Noventa y cinco años desde que se la fijó, y la categoría del municipio no cambió de nombre}, aunque hoy lo gobiernen un intendente electo y un Concejo Deliberante (capítulo \ref{cap:poblacion}).

\section{1943--1945: el departamento como línea y como resto}

Tres actos de policía, que no hablan de suelo ni de agua, dicen lo mismo que este capítulo por
otra vía.

\textbf{El Decreto 208-G, del 4 de agosto de 1943}, reorganiza las jurisdicciones de las
Comisarías Inspectoras de Zonas y divide la provincia en siete. Seis tienen inspección propia y
asiento en una localidad: Chicoana, Cachi, Rosario de la Frontera, y así. La Primera Zona no.
La integran La Capital, Campo Santo, La Caldera y Cerrillos, y el decreto dispone que «\textit{serán
atendidos directamente por Jefatura de Policía}». \textbf{El departamento no está mal atendido:
está atendido desde la capital, sin instancia intermedia.} Es la misma forma que tendrá,
setenta años después, la administración de su tierra fiscal, de su línea de ribera y de su
catastro.

\textbf{El Decreto 5738-G, del 10 de enero de 1945}, delimita el radio de las comisarías de la
Capital. Su inciso c) fija el \textbf{río Vaqueros} como límite norte del Destacamento Norte.
Un curso que nace y corre en La Caldera aparece en el Boletín, por primera vez en las ediciones leídas,
\textbf{como el borde de otra jurisdicción}: no como recurso, no como riesgo, sino como línea
que le sirve a la ciudad para saber dónde termina.

Y el mismo inciso excluye de esa jurisdicción a la \textbf{Destilería Chachapoyas}, «\textit{que
posee jurisdicción policial propia}». El decreto no dice quién la ejerce, bajo qué norma ni
desde cuándo: la da por existente. \textbf{Un enclave industrial con policía propia, sobre el
límite entre la capital y el departamento, en 1945.}

\begin{cautela}
Chachapoyas reaparece un mes después, ya integrada a la Primera Zona junto con La Caldera y
Vaqueros. El barrido no contiene el acto que explique el cambio, si lo hubo, y este libro no
estableció qué era exactamente esa destilería, de quién era ni qué relación tenía con el
territorio del departamento. \pendiente{antecedentes de la Destilería Chachapoyas y de la
jurisdicción policial propia que el Decreto 5738-G le reconoce}
\end{cautela}

Súmese el \textbf{Decreto 9686-G, del 15 de diciembre de 1945}, que convoca a las elecciones
del 24 de febrero de 1946 y reparte la representación en tres escalones. La Capital elige siete
diputados; cinco departamentos eligen dos; y dieciséis ---La Caldera entre ellos--- eligen
\textbf{uno}. El capítulo \ref{cap:politica} muestra que esa proporción sigue vigente ochenta
años después.

\section{1940: cuatro años antes, el departamento estaba en el plan}

Antes de leer el acto de 1944 conviene poner lo que había cuatro años antes, porque de otro
modo el capítulo estaría contando una decadencia que el archivo no respalda.

\begin{ficha}{Decreto Nº 4217 --- Salta, 12 de noviembre de 1940 --- aprueba el Acta Nº 376 de Vialidad de Salta}
El Acta, de la sesión del 15 de octubre de 1940 presidida por el ing.\ Eduardo Arias, aprueba
en su punto 14º un \textbf{plan integral de quince años} para la vigencia de la Ley nacional
12.625, por un total de \$7.477.000. Entre sus renglones:

\begin{itemize}[leftmargin=*,itemsep=1pt]
\item Salta--Yacones, 2º tramo, San Lorenzo--Castellanos: \textbf{\$10.000}
\item Salta--Yacones, 2º tramo, Castellanos--Yacones: \textbf{\$90.000}
\item \textbf{Puente Río Caldera: \$200.000}
\item Puente Río Mojotoro: \$260.000
\end{itemize}

El punto 15º aprueba el plan de Coparticipación Federal para 1941, de \$714.000, que incluye
los \$90.000 del tramo Castellanos--Yacones. Y los puntos 4º y 5º liquidan certificados
parciales de dos obras \textbf{en ejecución}: el camino de Salta a Los Yacones, tramo Salta--San
Lorenzo, al contratista Antonio Rosselló; y el \textbf{camino de Calderilla a Río Las Pavas},
tramo Desmonte a Campo Santo, a los contratistas Cvitanic y Herceg, por \$12.807,74.

El punto 16º aprueba el convenio de transferencia de caminos con Vialidad Nacional. Tras él,
la \textbf{Ruta Nacional Nº 55} cruza «\textit{Cerrillos, Salta, Vaqueros, Calderilla, hasta
Tres Cruces en el límite con Jujuy}», y la red provincial conserva «\textit{La Calderilla a Río
Saladillo por Gallinato y Campo Santo}», «\textit{Salta a Los Yacones por San Lorenzo}»,
«\textit{Puente Vaqueros a Castellanos}» y el «\textit{Acceso a Estación Mojotoro}».
\end{ficha}

\textbf{Y venía de cinco años de obra.} El camino de La Calderilla a El Desmonte por Gallinato
aparece en dieciséis entradas de las actas 100 a 110 del Directorio de Vialidad, publicadas
entre septiembre y diciembre de 1935: jornales de cuadrilla, alambrado, portland para muros, presupuesto de alcantarillas, un
puente de veinte metros de luz, sobrestantes, una licitación desierta y la adjudicación de las
obras de arte. Se estaba construyendo, semana a semana, con planilla.

\textbf{Y tiene muertos.} En noviembre de 1935 la Provincia transfiere a la Caja Nacional
\textbf{\$2.500} de indemnización por el accidente del que fue víctima el obrero \textbf{José
Gallardo} en el camino Gallinato, conforme la Ley 9.688. En marzo de 1936 el Directorio paga
\$100 a la Empresa Guzmán «\textit{por tres ataúdes para los \textbf{tres obreros fallecidos en
el accidente de Gallinato}} y un servicio fúnebre a uno de los mismos», y en la misma acta
designa un capataz «\textit{en reemplazo de don Juan Di Bes reciente accidentado en el camino
por Gallinato}» ---el mismo capataz cuya cuadrilla figura dos ediciones antes trabajando en el
camino de La Calderilla. \textbf{Son las únicas muertes obreras que este relevamiento encontró
en todo el archivo}, y aparecen por la factura de los ataúdes.

\begin{cautela}
Los avisos no dicen en qué departamento ocurrieron los hechos: el camino por Gallinato es la
traza histórica de Salta hacia General Güemes y cruza el límite con Campo Santo. Tampoco dicen
si el accidente de Gallardo en 1935 y el de los tres obreros en 1936 son el mismo hecho o dos
distintos: las piezas están en actas consecutivas del mismo Directorio y el corpus no los
vincula. Este libro no los suma.
\end{cautela}

\textbf{En 1940 el departamento no era el resto de nada.} Tenía dos obras en ejecución con
certificados parciales pagados, cuatro caminos en la red provincial, una ruta nacional que lo
atraviesa de punta a punta ---la que hoy se llama la Cornisa--- y \textbf{doscientos mil pesos
asignados a un puente sobre su río} dentro de un plan a quince años.

Los boletines de 1941 y 1942 siguen la ejecución paraje por paraje: jornales del camino de
\textbf{La Caldera al Sauce}, transporte de ripio y piedras del camino de \textbf{Vaqueros a
Lesser}, un peón caminero para ese tramo, jornales del camino de \textbf{Lesser a Los Yacones} y
del de \textbf{Yacones a Potrero de Castillo}, cambio de contratistas en el de Calderilla a Río
Las Pavas. No era un plan en el papel: era obra con planilla de jornales.

\begin{cautela}
Tres reservas. La columna de importes del plan de quince años sale desalineada en la capa de
texto del Boletín y los cinco valores citados se leyeron sobre la imagen; la suma del plan de
1941 sí cierra (255.000 + 90.000 + 200.000 + 50.000 + 119.000 = 714.000), lo que respalda esa
segunda tabla, pero el total del plan de quince años no fue recalculado renglón por renglón.
Segunda: varios de estos caminos ---Salta a San Lorenzo, por ejemplo--- corren fuera del
departamento aunque su destino esté adentro. Y tercera: \textbf{este libro no estableció qué se
construyó}. Que una obra esté en un plan y tenga certificados parciales no prueba que se haya
terminado, y el puente de \$200.000 sobre el río La Caldera no vuelve a aparecer en las
mil cuatrocientas setenta y dos ediciones relevadas. \pendiente{ejecución efectiva del plan de
quince años de la Ley 12.625 en el departamento, y destino de la partida del Puente Río
Caldera}
\end{cautela}

\textbf{Y eso es lo que le da su peso al acto de 1944.} Los vecinos habían pedido un
puente en 1935 y la respuesta fue «se dispondrá en su oportunidad». En 1940 la oportunidad
llegó al papel. En 1944 se retira del plan el camino de Vaqueros a Los Yacones. \textbf{No es
una decadencia lenta: es una reversión, y tiene fecha y fundamento escrito.}

\section{1944: la Provincia explica por qué}

Hay un acto que dice, con una franqueza que no se repite en todo el archivo, lo que este
capítulo viene argumentando.

\begin{ficha}{Decreto Nº 3016-H --- Salta, 4 de mayo de 1944 --- Expte.\ Nº 15.993/1944}
Aprueba la resolución del punto 5º del Acta Nº 31 del Consejo de la Administración de Vialidad
de Salta, del 25 de abril de 1944, que transcripta dice:

«\textit{CAMINO DE VAQUEROS A LOS YACONES. Atendiendo el informe producido por la División
Estudios y Proyectos que establece que la construcción de este tramo prevista en su oportunidad
para el Plan del corriente año de la Ley 12.776, \textbf{demandaría un gasto sumamente elevado
por las características del terreno, gasto que no estaría justificado con el beneficio que la
inversión reportaría}, se resuelve eliminarlo de las obras de dicho plan y afectar el importe
de \$40.000.--- m/n., a la construcción del camino de Salta a La Floresta, que de este modo
reemplazaría al anterior.}»

Firman el Gral.\ José Morales Bustamante y Carlos A.\ Emery. Aparición única, B.O.\ Nº 2.054
del 12 de mayo de 1944.
\end{ficha}

Cuarenta mil pesos de obra vial, ya presupuestados y ya asignados al plan de 1944, se retiran
del camino que uniría dos localidades del departamento y se pasan a una obra que está fuera de
él. \textbf{El departamento perdió su camino, y el expediente dice por qué}: porque la
inversión no se justificaba con el beneficio que reportaría.

Tres cosas conviene retener del acto, y ninguna es la decisión en sí.

\textbf{Primera.} El criterio se enuncia. No hay aquí opacidad: la Provincia escribe que aplicó
una relación costo-beneficio y que el departamento la perdió. \textbf{No es la única vez en este
libro que un organismo publica el criterio por el cual La Caldera queda afuera de algo}: en 1918 el
decreto 1.660 publicó la cuenta de renta por la que el departamento no podía tener municipalidad
electiva, y en 1919 el decreto 33, el umbral de población que dejaba fuera a todos los
departamentos de campaña (capítulo \ref{cap:opacidad}). Las tres veces el criterio de la exclusión
se escribe en el acto que excluye, que es lo contrario de lo que el capítulo \ref{cap:opacidad}
describe para 2014--2026. Una versión anterior de este párrafo decía que el de 1944 era el único.

\textbf{Segunda.} El Boletín no dice si el mismo criterio se aplicó a la obra beneficiaria. El
camino de Salta a La Floresta recibió los fondos sin que conste que se le haya calculado nada.

\textbf{Tercera.} La plata era nacional ---Ley 12.776--- y estaba afectada a un plan aprobado.
La resolución manda \textbf{comunicar la modificación al Ministerio de Obras Públicas de la
Nación} una vez aprobada por el Ejecutivo provincial: se reasigna primero y se avisa después.

El proyecto no vuelve a aparecer en las mil cuatrocientas setenta y dos ediciones del barrido. Y
ochenta años después, el diagnóstico social del Plan de Manejo registra que los tres
puentes del departamento son «antiguos, obsoletos, sin mantenimiento y muy estrechos».

\section{1940 y 1942: también las truchas}

Queda una pieza chica, de doscientos pesos, que dice lo mismo que todo lo anterior y además
nombra a quiénes.

En julio de 1940 la Provincia autoriza \$673,45 para traer desde Córdoba al \textbf{Dr.\ Tomás
L.\ Marini}, Jefe de la División de Piscicultura del Ministerio de Agricultura de la Nación, y
para comprar diez mil truchas. Su informe, transcripto íntegro en el decreto, enumera dónde ya
hay pejerrey en la provincia ---«\textit{en las fincas de los señores: Dr.\ Patrón Costas, en
Uratay; \textbf{Lucio Ortiz, en Yacones}; Indalecio Gómez, en Pampa Grande}»--- y qué
instalaciones hay disponibles para mantener truchas en cautividad. La tercera de las cuatro es
el «\textit{\textbf{canal de propiedad del Dr.\ Lucio Ortiz en Yacones}}».

Dos años después, el \textbf{Decreto 6749-H del 23 de noviembre de 1942} paga \$200 al
Ministerio de Agricultura de la Nación por «\textit{los 5.000 alevinos de trucha arco-iris que
fueron sembradas en esta Provincia \textbf{en los ríos Castellanos y Wierna}}».

Tres cosas se leen ahí.

\textbf{Los ríos del departamento se poblaron con salmónidos exóticos en 1942}, por decisión
nacional y con pago provincial, y el decreto lo autoriza \textit{después} del hecho: habla de
alevinos «que fueron sembradas», en pasado. No dice en qué tramo, en qué fecha ni si hubo
estudio previo.

\textbf{Las aguas que el Estado se disponía a poblar eran privadas, y sus dueños eran el
gobierno.} El informe nombra a un ex gobernador, a un ministro y a un tercer doctor con canal
propio en Yacones. El decreto que lo aprueba lo refrenda \textbf{Jaime Indalecio Gómez}, y en
la lista de beneficiarios figura un \textbf{Indalecio Gómez} en Pampa Grande. Como siempre en
este libro: la coincidencia de apellido se anota, la identidad no se afirma, y se resuelve con
una partida.

\textbf{Y el canal del Dr.\ Lucio Ortiz en Yacones queda documentado como propiedad privada con
obra hidráulica propia en 1940}, en la subcuenca que la presentación del Plan Maestro de Drenaje
llamará Volcán--Yacones.

\section{La conclusión que se sigue}

El problema de La Caldera no es que el Estado esté ausente: el capítulo \ref{cap:opacidad} muestra una relación densa, continua y antigua.

Tampoco es que el municipio sea incapaz, en el sentido de que le falte voluntad o
competencia. Tiene código de ordenamiento, plan urbano, cuadrilla vial, centro de emisión de
licencias y capacidad para ejecutar obra pública por contrato.

\textbf{El problema es que el municipio nunca fue diseñado para decidir, y el territorio
nunca fue tratado como un lugar donde vive gente.} Fue tratado, desde que hay registro
escrito, como el lugar de donde sale algo que otros necesitan: agua, piedra, aire de fin de
semana, suelo barato y cercano.

Eso explica cosas que de otro modo parecen contradictorias. Explica que el Estado provincial
pueda gastar dos millones de pesos en proteger de las crecidas las cámaras de captación del
acueducto en 2014, y que cuatro años después el río se lleve casas: lo que se protegía era
la infraestructura, no la población. Explica que exista una autoridad metropolitana con
planificación, regulación y poder de policía para el transporte de pasajeros, y ninguna para
la cuenca: los colectivos llevan gente a la capital; la cuenca es un problema local.
Explica que la primera concesión provincial de agua que el archivo registra en el departamento, en 1924, se otorgara con el
argumento de que «permitirá fácilmente colonizar esas tierras»: el agua se concedía para que
la tierra pudiera ocuparse, no porque hubiera alguien a quien darle agua.

Y explica, finalmente, por qué las dos velocidades del capítulo anterior conviven sin
escándalo. Quien compra un lote de seiscientos metros para pasar los veranos es, desde el
punto de vista del dispositivo, el uso previsto del territorio. Quien vive ahí todo el año y
espera treinta años una escritura, no es el destinatario de nada: es el residuo de un lugar
que se administra para otra cosa.

\section{Un vecino lo dijo mejor}

Mientras este libro se escribía, alguien formuló la misma idea en cuatro frases y desde
adentro.

El 27 de enero de 2026, en plena emergencia por las crecidas, el presidente del Centro
Vecinal de La Calderilla hizo declaraciones a una radio de Salta que la prensa recogió. Son
testimonio y no documento, y como tal se las presenta: no acreditan hechos, acreditan que
alguien que vive ahí lee su situación de esta manera.

\begin{ficha}{Declaraciones del presidente del Centro Vecinal de La Calderilla, 27 de enero de 2026}
Sobre la aprobación de los loteos: «Hay un montón de gente ofreciendo loteos, pero
\textbf{el municipio no nos contesta si están aprobados o no}. Y si no lo están, la pregunta
es por qué no los clausuran.»

Sobre los servicios: «\textbf{Nosotros no tenemos agua} y, sin embargo, \textbf{se ofrecen
loteos con redes de agua, cloacas, arena del Paraná y hasta palmeras}. Es una locura.»

Sobre el monte ribereño: «El monte ribereño no se puede desmontar porque protege frente a
una crecida que puede ocurrir cada 15, 20 o 50 años. Estos emprendimientos desmontaron todo
el lote, incluso esa franja, \textbf{sin autorización de la Secretaría de Medio Ambiente}.»
Y agrega que, pese a que el emprendimiento habría sido multado por el organismo ambiental,
el municipio avanzó igual con una audiencia pública, en la que «la presidenta del Concejo
Deliberante \textbf{agradeció la inversión}, cuando el proyecto todavía estaba penalizado y
no había recuperado el monte ribereño».

Sobre el agua de la capital: «\textbf{La planta potabilizadora está en La Caldera} y también
hay \textbf{cavernas subterráneas debajo del lecho del río que captan agua por filtración}.
\textbf{Estos loteos están aguas arriba de esas tomas}.» Y cierra: «Esto no es sólo un
problema de La Caldera. Es importante que la gente de Salta capital se entere, porque está
en juego la reserva de agua potable.»

\textit{La nota consigna el apellido del dirigente de dos maneras distintas --- Barrios en el
copete, Barros en el cuerpo ---. La grafía correcta es la segunda: el Ente Regulador de los
Servicios Públicos lo identifica, en enero de 2025, como el Ing.\ \textbf{Miguel Ángel Barros},
presidente del Centro Vecinal La Calderilla, en una reunión con concejales del departamento en
la que dijo haber salido informado «sobre el tema de los desarrolladores de los loteos».}
\end{ficha}

\begin{cautela}
\textbf{La audiencia no está en el Boletín.} El buscador de avisos del Boletín Oficial provincial no devuelve ninguna convocatoria a audiencia pública en La Caldera publicada entre 2023 y 2026 (consulta de septiembre de 2026). Fuera de los tramos leídos edición por edición, una ausencia sólo significa que el barrido no la encontró; pero en este caso deja tres posibilidades, y este libro no puede elegir entre ellas: que la audiencia haya sido convocada por el propio municipio sin publicación provincial; que sea anterior a 2023; o que el dirigente llame «audiencia pública» a otra instancia, como una sesión del Concejo. Tampoco puede identificar a la presidenta del Concejo que menciona, pero puede acotar cuándo hubo una. La prensa oficial registra, en este orden: a \textbf{Yamila Sastre} como presidenta en diciembre de 2021 (portal de la Municipalidad de La Caldera); a \textbf{Ana Sumbay} como presidenta, en una nota del Ministerio de Desarrollo Social que no permite fechar su mandato; a \textbf{Germán Sumbay} como presidente en enero de 2025 (Ente Regulador) y en enero de 2026 (capítulo \ref{cap:expropiacion}); y a \textbf{Jorge García} como presidente en agosto de 2026 (Escribanía de Gobierno). Si el relato es exacto, la audiencia es anterior a 2025. Y el buscador de avisos, extendido hasta 2019, devuelve \textbf{una sola convocatoria a audiencia pública sobre un loteo del departamento}: la de Lares de la Inmaculada, en Vaqueros, de marzo de 2019 (capítulo \ref{cap:vaqueros}), que no corresponde ni al municipio ni al paraje del relato. La otra del período, la de Vialidad por la Ruta 9 de noviembre de 2023, también se celebró en Vaqueros y trató una obra vial. Lo más probable, entonces, es que la instancia que el dirigente llama «audiencia pública» haya sido municipal ---una sesión o una reunión convocada por el Concejo--- y no la audiencia ambiental provincial de la Ley 7070, que se publica en el Boletín. \pendiente{fecha y carácter de la instancia}
\end{cautela}

Las cuatro frases corresponden, una por una, a las cuatro cosas que este libro tardó más de quinientas páginas en documentar.

La primera es la \textbf{segunda tesis}: la opacidad no está en la existencia del acto sino
en la imposibilidad de saber si existe. Un vecino que pregunta al municipio si un loteo está
aprobado y no recibe respuesta está describiendo, sin nombrarla, la misma carencia que este
trabajo encontró en seis instrumentos con anexos sin publicar.

La segunda es la \textbf{primera tesis} en su forma más cruda. Quien vive ahí desde siempre
no tiene agua; quien compra un lote nuevo la recibe, con cloacas y arena traída del Paraná.
No es una metáfora del libro: es la descripción literal de los dos circuitos, hecha por
alguien que está de un lado de la línea.

La tercera toca el punto donde este libro es más cauto y donde ese vecino no lo es. Sostiene
que se desmontó monte ribereño sin autorización, que hubo multa ambiental y que el municipio
llevó igual el proyecto a audiencia pública. \textbf{Este libro no verificó nada de eso} ---
no obtuvo el acta de la multa, ni el expediente, ni la versión taquigráfica de la audiencia
--- y por lo tanto no lo afirma. Lo consigna porque señala tres documentos concretos que
existen y que resolverían el punto, y porque, de confirmarse, describiría el mecanismo
completo: sancionar y autorizar a la vez.

Y la cuarta es \textbf{la tesis de este capítulo}, dicha sin rodeos por alguien que no
escribió un libro para llegar a ella. La planta potabilizadora está ahí. Las tomas por
filtración están debajo del lecho del río. Los loteos están aguas arriba de las tomas. Y la
conclusión que ese vecino extrae es exactamente la que se extrae de más de quinientas páginas de expedientes: \textbf{esto no es sólo un problema de La Caldera}.

\section{Y hubo un momento en que no era así}

Todo lo anterior describe un estado de cosas, y podría leerse como si hubiera sido siempre
el mismo. No lo fue, y el archivo permite fechar la diferencia.

En 1914 el municipio de La Caldera reglamentó por ordenanza el reparto \textbf{íntegro} del
caudal de uno de sus ríos, reconociendo los usos consuetudinarios que ya existían. En 1939
esa ordenanza le ganó a una solicitud de concesión provincial: el Poder Ejecutivo denegó el
pedido de un comerciante de la capital porque otorgarlo habría alterado el reparto local, y
al año siguiente ratificó que \textbf{no podía} conceder agua de ese río.

En 1924 y en 1938, el dictamen de la comisión municipal integró el expediente de dos
concesiones de agua y quedó publicado. En 1939, un certificado \textbf{firmado por Augusto
Regis}\index{Regis, Augusto}, presidente de la Comisión Municipal, fue el instrumento que un
propietario presentó ante la Provincia para acreditar de dónde tomaba el agua su finca: el
barrido de ese año lo ubica en el Boletín Nº 1778, hoja 15, y transcribe que el certificado
acredita «\textit{la principal provisión y suministro}» (capítulo \ref{cap:siglo}).

\textbf{Ese municipio hablaba sobre el agua de su territorio, y lo que decía pesaba.}

Noventa años después no aparece en ninguno de los veinticinco actos de línea de ribera que la Provincia dictó sobre su cuenca entre 2013 y 2015 ---salvo en la Resolución 023/14, que la serie cuenta aparte y que el intendente pidió para el propio pueblo---; el plano del loteo más discutido de este
libro se aprobó por resolución ministerial sin ninguno de los diez actos que su
propio Código exige; y un vecino declara en enero de 2026 que pregunta al municipio si los
loteos están aprobados y no le contestan.

De modo que la frase que da título a este capítulo necesita una fecha. \textbf{La Caldera no
fue siempre infraestructura de una ciudad que no la gobierna}: fue también, durante al menos
un cuarto de siglo, un municipio pobre y tutelado que sin embargo decidía sobre el agua que
corría por su suelo.

\begin{cautela}
\textbf{Y ese cuarto de siglo tiene un agujero en el medio, que el barrido de 1922 abrió al
cierre y que obliga a no leer esta sección como una curva descendente.} La ordenanza que reparte
el caudal íntegro del río Santa Rufina es de octubre de 1914. \textbf{Ocho años después, el 18 de
abril de 1922, un Interventor Nacional tiene que crear la Comisión Municipal del departamento
porque «\textit{en la actualidad carece de ella}»} (capítulo \ref{cap:siglo}). El cuerpo que
había dictado aquella ordenanza no estaba disminuido: no existía, y ningún acto publicado dice
cuándo dejó de existir.

\textbf{Tres meses más tarde ese cuerpo, recién nombrado desde afuera, sesiona, sanciona su
presupuesto y se lo aprueba el Poder Ejecutivo sin publicarlo. Y diecisiete meses después la
Provincia lo interviene}: el decreto 1161 del 24 de septiembre de 1923 declara intervenida la
Municipalidad del departamento y le nombra un Interventor, sin publicar ni el expediente que lo
funda ni las instrucciones que le da ni un solo acto suyo.

De modo que entre el municipio
que en 1914 repartía un río y el que en 1939 certificaba de dónde tomaba el agua una finca hay,
al menos, un año sin municipio y otro con el municipio intervenido. \textbf{Lo que el archivo muestra no es una
pérdida progresiva de voz sino una presencia intermitente}, y eso cambia la pregunta: no es sólo
cuándo dejó de decidir, sino cuántas veces dejó de existir.

\textbf{Y hay un caso donde las dos cosas se tocan.} El informe municipal que fundó la primera
concesión de agua del departamento ---el dictamen favorable que el decreto de enero de 1924
invoca--- se emitió en el lapso en que esa comuna estaba intervenida por la misma Provincia que
concedía, y el archivo no dice si lo firmó la Comisión o el Interventor (capítulo
\ref{cap:siglo}). \textbf{La voz municipal sobre el agua, en su primer ejercicio documentado ante
la Provincia, puede no haber sido municipal.}


\textbf{Todo ese arco cabe en la lámina \ref{fig:agua}}, que pone en un mapa las fincas de cada
acto y, en una línea de tiempo, los seis momentos con su verbo. \textbf{Lo que se ve de un golpe
es que el punto donde la facultad cambia de manos ---las juntas del Wierna con el río La
Caldera--- está dentro de este departamento}, y que las fincas que el municipio regó entre 1885
y 1938 están todas aguas arriba de él.

\textbf{Y la ordenanza de 1914 seguía viva veinticinco años después.} En 1939, al
resolverse la concesión de agua para regadío de la finca «Campo Alegre», la
oposición al pedido invoca «\textit{la existencia de una ordenanza dictada por el
municipio de la Caldera en fecha octubre 13 de 1914, que reglamenta la
distribución de la totalidad}» del caudal, y la Provincia tiene que pronunciarse
sobre ella (capítulo \ref{cap:siglo}). \textbf{No es sólo que el municipio
dictó aquella ordenanza: es que un cuarto de siglo después todavía se la
esgrimía como derecho vigente en un expediente provincial, y contra una
concesión nueva.} Es la prueba más fuerte que este libro tiene de que aquella
voz municipal no fue un gesto aislado.

\textbf{Y tres meses después el beneficiario de esa concesión entró al cuerpo que la había
informado.} El decreto 1.576 del 15 de abril de 1924 designa a Francisco Urquiza\index{Urquiza, Francisco S.}
miembro de la Comisión Municipal de La Caldera. No estaba adentro cuando el cuerpo informó ---el
decreto que concede es del 28 de enero--- y sí estaba adentro cuatro años más tarde, cuando esa
misma comisión volvió a informar a su favor en el reclamo de Campo Santo (capítulo
\ref{cap:politica}). \textbf{De modo que el único período en que este libro puede mostrar al
municipio decidiendo sobre el agua de su suelo es también aquel en el que el cuerpo que decide y
el titular del derecho se superponen.} Eso no invalida la sección: la afila. \textbf{Lo que el
archivo prueba no es que hubiera una voz municipal autónoma, sino que había una voz municipal}, y
que lo que se perdió después fue incluso eso.

\begin{cautela}
\textbf{Y el tramo 1937--1946 agrega la parte que a esta sección le faltaba: qué hizo el Estado
con el río, además de con el agua.} \textbf{Tres veces se proyectó defenderlo y ninguna obra
consta terminada.} En 1938, a pedido de vecinos y porque «\textit{todos los años en la época de
lluvias las aguas del río Caldera amenazan las propiedades}», se proyectan \textbf{cincuenta
metros} de defensa sobre la margen izquierda. En 1940 el propio Estado informa que aquella obra
«\textit{fue suspendida en virtud de que el mencionado río desvió su curso}» y autoriza un
\textbf{reparo provisorio} «\textit{en el lugar denominado la Calderilla}». En 1945 hay
licitación pública por \textbf{\$2.515} y la obra termina ejecutándose «\textit{por vía
administrativa}», sin certificado ni recepción (capítulo \ref{cap:siglo}).

\textbf{Eso cambia la forma de la pregunta que abre el capítulo \ref{cap:amparo}.} Ese capítulo
parte del aluvión de diciembre de 2018 y del Plan de 2024 que cifra la defensa del río en el
cincuenta y ocho por ciento de veintitrés mil millones. \textbf{Lo que el archivo muestra es que
el problema no empezó en 2018 y que el Estado lo reconoció por escrito en 1938}: ochenta años,
tres proyectos y ninguna obra que conste terminada. \textbf{Y la escala de aquella decisión está
en el mismo Boletín}: la defensa del río valía \$2.515 cuando refaccionar la comisaría del
pueblo, el año anterior, había costado \$5.573,80 con su adicional (capítulo \ref{cap:defensas}).
\end{cautela}

\textbf{Y en el mismo año aparece la otra mitad del reparto.} El decreto 1.580 del 16 de abril de
1924 presta la aquiescencia para instalar en \textbf{San Alejo} una escuela de la Ley 4874 ---la
ley Láinez, de escuelas nacionales en territorio de provincia---. \textbf{La tercera escuela del
departamento no la pone la Provincia que se lleva su agua: la pone la Nación}, y la Provincia se
limita a consentir. Es la primera vez que este relevamiento ve al Estado nacional poniendo
infraestructura dentro del departamento, y es una escuela.

\begin{cautela}
\textbf{Y en 1929 la Legislatura escribió la tesis de este capítulo en el texto de una ley.} La
\textbf{Ley 11.078}, del 27 de septiembre de 1929, autoriza a la \textbf{Municipalidad del
Departamento de Campo Santo} a hacer los trabajos que conceptúe necesarios en la playa del río
Mojotoro «\textit{para reconcentrar el agua, \textbf{desde las juntas del Río de Wierna con el de
La Caldera}}», y a «\textit{vigilar y prohibir se abran nuevas bocas-tomas}» (capítulo
\ref{cap:siglo}).

\textbf{El punto de arranque está dentro de La Caldera y la facultad es de otro municipio.} No
hay que interpretar nada: la ley le da a Campo Santo el poder de policía sobre las tomas de un
agua que nace y baja por este departamento, y a este departamento no le da ninguno. \textbf{Es el
documento más antiguo y más explícito que este libro tiene de la asimetría que describe}, y es de
un año en que La Caldera todavía tenía comisión municipal en funciones.

\textbf{Y en 1931 el arreglo se vuelve institución}: una Comisión \textit{ad hoc} sobre las aguas
del Vaqueros y del Wierna, integrada por los Comisionados Municipales de los dos departamentos
---los dos designados, ninguno electo---. \textbf{De modo que la pregunta de esta sección tiene
ahora una respuesta parcial con fecha}: la voz municipal sobre el agua no se apagó de golpe, se
fue repartiendo, y en 1929 la parte que se repartió quedó del lado de abajo del río.
\end{cautela}


\end{cautela}

Cuándo dejó de hacerlo es la pregunta que este libro deja abierta, y no por descuido: ocurre
en el período que su relevamiento no cubre. Entre la ordenanza de 1914 y los actos de 2013 hay casi un siglo, y de él este relevamiento dejó sin leer de corrido los sesenta y dos años que van de 1951 a 2012 ---1947 a 1950 se barrieron al cierre, sin lectura sobre la imagen, con la reserva que el apéndice \ref{cap:fuentes} declara---.

\section{Un sector donde se ve todo junto: El Gallinato}

La tesis de este capítulo se puede leer de una sola vez en un sector del departamento. La
fracción censal 66077-03 ---el sector oriental, que contiene El Gallinato y su escuela (capítulo
\ref{cap:ausencias})--- reúne cuatro piezas que el libro documentó por separado.

\textbf{El camino.} Desde los años treinta y hasta 1942, el Directorio de Vialidad construyó
acta por acta el camino de La Calderilla a El Desmonte y a Campo Santo por Gallinato, con
puentes, alcantarillas y cuadrillas; en él murieron obreros, aunque los avisos no dicen de qué
lado del límite departamental (capítulo \ref{cap:presencia} y apartado de 1940 de este capítulo).
\textbf{El camino no termina en Gallinato}: sigue hacia Campo Santo y General Güemes.

\textbf{La escuela.} En 1939 la escuela para Gallinato la tramitó la Nación, por la Ley Láinez, y
la Provincia se limitó a prestar la aquiescencia (capítulo \ref{cap:presencia}). La escuela que hoy funciona en el
paraje, la Nº 4118, estuvo desde 1968 en terreno ajeno, y la Provincia adquirió ese suelo por
prescripción recién en 2015 (capítulo \ref{cap:tierrafiscal}).

\textbf{El agua.} Es la fracción con peor cobertura de red del departamento: el 58 por ciento de
sus hogares, sin conexión (capítulo \ref{cap:aguabaja}).

\textbf{Y la gente.} Tiene 764 habitantes en viviendas particulares, entre ocho y diez puntos
menos de población de 15 a 29 años que las otras dos fracciones, la mayor proporción de mayores de
65 y el 43,17 por ciento de sus viviendas sin residentes permanentes (capítulo \ref{cap:ausencias}).

\begin{cautela}
\textbf{Pasó el camino y no el agua.} Es la forma que este capítulo le atribuye al departamento
entero ---línea para lo que va a otra parte, resto para lo que se queda---, vista en un sector que
se puede recorrer en una tarde.

Tres reservas. La fracción contiene El Gallinato, pero no es sólo El Gallinato. En un universo de
764 personas, un puñado de casos mueve varios puntos porcentuales. Y cada pieza tiene su fuente,
pero el orden en que se las pone, y lo que ese orden sugiere, es lectura del autor: \textbf{nada de
esto prueba que alguien haya decidido que el camino pasara y el agua no}.
\end{cautela}

\section{Quiénes pagan}

Son doce mil doscientos noventa y nueve. Los que se inundaron dos veces en trece días, el 21 de diciembre de 2018 y el 3 de enero de 2019, y otra vez en enero de 2026. Los que llevan más de quince años reclamando la
red de agua en la zona alta de Vaqueros mientras un club de campo del mismo municipio obtuvo
del Estado la exigencia —no el permiso: la exigencia— de resolver su abastecimiento en forma
privada. Los que llevan más de siete años y medio litigando en una vía que la Constitución define
como expedita. Los caldereños que, según su centro vecinal, ocupaban tierra fiscal en un barrio de la capital y obtuvieron en 2025 la adjudicación en venta, en cien cuotas y con hipoteca ---en un caso, tras treinta años de expediente---, después de acreditar que no tenían otro inmueble. Y los que, sin ocupar, piden un lote por sorteo y deben acreditar residencia policial,
declarar bajo pena de falso testimonio y dar referencias vecinales con nombre, apellido y
franja horaria de visita domiciliaria.

Ninguno de ellos aparece en las fuentes de este libro salvo cuando reclama, litiga o
solicita. Es la carencia que el capítulo \ref{cap:metodo} declara y que el capítulo \ref{cap:resistencias} desarrolla: no hay
aquí entrevistas propias, sólo las dos conversaciones informales del capítulo \ref{cap:poblacion}.

\section{Una empresa que cierra el circuito}

Este capítulo sostiene que el departamento funciona como insumo de la capital. Hay una firma en
la que ese circuito se cierra entero, y los tres tramos están documentados en el Boletín
Oficial.

\begin{ficha}{\textsc{transporte hnos.\ mendoza s.r.l.} --- tres roles documentados}
\textbf{Uno.} Tramitó ante el Juzgado de Minas la \textbf{renovación de la concesión de la
cantera de áridos «Macarena»}, en el departamento La Caldera, \textbf{lugar río Wierna}, sobre
diez hectáreas fiscales \textit{(cap.\ \ref{cap:resistencias}; el cuadro minero de 2018 la nombra «Hermanos Mendoza S.R.L.», y que sean la misma sociedad es probable y no está verificado)}.

\textbf{Dos.} Su razón social es la de una \textbf{empresa de transporte}
\textit{(cap.\ \ref{cap:ausencias})}.

\textbf{Tres.} El \textbf{Boletín Oficial Municipal de la ciudad de Salta Nº 2.725}, del 27 de
agosto de 2025, publica la Resolución 0207 de la Secretaría de Obras Públicas, que le otorga
una ampliación de plazo de veintiún días \textbf{para terminar la obra «bacheo con hormigón y
bacheo con mezcla asfáltica con polímeros de aplicación en frío en zona C.S.O.»} de la ciudad
de Salta.
\end{ficha}

\textbf{Y no es una obra sola.} El buscador del Boletín Oficial Municipal devuelve la razón
social en \textbf{ocho ediciones entre 2022 y 2026}. Dos fueron verificadas en su texto: la
Resolución 0207 de agosto de 2025 ya citada, y la \textbf{Resolución 0143 del 2 de septiembre de
2026}, que aprueba la planilla de economías y demasías de la obra «\textbf{puesta en valor
Plaza Dr.\ Bernardo Frías}», \textbf{adjudicada a la misma firma por orden de compra 003/2025}.

\textbf{Extrae el árido del cauce en La Caldera, lo transporta, y ejecuta obra pública en la
capital de manera sostenida durante al menos cuatro años.} Es la tesis de este capítulo reducida
a una sola firma.

\subsection*{Y no es una firma sola}

Buscada en el mismo registro, \textbf{la segunda concesionaria de áridos del corredor aparece
también como contratista de obra pública de la capital.}

\begin{ficha}{Boletín Oficial Municipal de la ciudad de Salta Nº 2.646, 2 de octubre de 2024}
Resolución 0110 de la Subsecretaría de Contrataciones: «\textbf{aprobar la addenda de contrato
de obra pública} celebrado entre la Municipalidad de Salta \ldots\ y la \textbf{Empresa INCOVI},
representada por el Sr.\ Héctor Juan Pablo Balderrama».
\end{ficha}

\textbf{INCOVI S.R.L.\ es, según el capítulo \ref{cap:resistencias}, concesionaria de dos
canteras sobre tierra fiscal} en este corredor: «El Vaquero», en Vaqueros, sobre nueve
hectáreas fiscales, y «Los Yacones», \textbf{sobre el cauce del río Wierna}, con veintiuna hectáreas y 6.374 metros cuadrados, aunque el informe de Minería de 2023 registra que la renovación de esta última fue rechazada en junio de 2020 (capítulo \ref{cap:amparo}). \textbf{Más de treinta hectáreas fiscales entre ambas.}

\textbf{Dos de las concesionarias de áridos del corredor ---las dos que este relevamiento ubica con canteras sobre el Wierna, ambas en la nómina de 2022---
son contratistas de obra pública de la capital.} La coincidencia deja de ser un caso.

\begin{cautela}
\textbf{Y acá vale la misma advertencia, reforzada.} No está probado que el material de esas
canteras sea el que se usa en esas obras, ni hay irregularidad en que una empresa de áridos
ejecute obra vial: es la integración esperable del rubro.

\textbf{Lo que el patrón muestra es la dirección del circuito.} Las canteras están sobre cauces
y tierra fiscal del departamento La Caldera y de Vaqueros; las obras, en la capital. \textbf{El canon lo percibe la autoridad minera provincial, salvo el convenio de control y cobro con el municipio que menciona la gacetilla de 2018} ---capítulo \ref{cap:resistencias}; el apéndice \ref{cap:pedidos} pide ese convenio---, \textbf{la obra la paga y la
recibe otro municipio}, y el departamento que aporta el material tiene en esa rama cincuenta y
una personas ocupadas.

\textbf{Una pista que se siguió y no llevó a ninguna parte, y que se consigna porque este libro
también registra eso.} En el mismo boletín de octubre de 2024, la Dirección General de
Contrataciones de Obras Públicas aprueba el pliego para el «\textit{alquiler de topadora sobre
oruga para encauzamiento del arroyo Isasmendi}», y el capítulo \ref{cap:ribera} registra un
arroyo Isasmendi en la \textbf{finca Wierna}, catastro 2020, dentro de la serie de líneas de
ribera del departamento.

\textbf{No son el mismo curso.} Un segundo instrumento de la misma serie ---la resolución que
adjudica la obra--- consigna la \textbf{«Ubicación: Barrio El Progreso, Zona Sur de la Ciudad de
Salta»}. Es un homónimo del ejido capitalino, en el extremo opuesto del área metropolitana.

\end{cautela}

\begin{cautela}
\textbf{Sobre las ocho ediciones.} El buscador devuelve el boletín completo, no el acto, de modo
que \textbf{ocho ediciones no equivalen a ocho contratos}: una obra puede generar varias
resoluciones ---adjudicación, ampliación de plazo, redeterminación, economías y demasías--- y
aparecer en varias ediciones. \textbf{Lo que queda establecido son dos obras distintas},
verificadas en su texto, y una presencia continuada en el registro entre 2022 y 2026.
\end{cautela}

\begin{cautela}
\textbf{Qué prueba esto y qué no, porque la tentación de estirarlo es fuerte.}

\textbf{Prueba} que la misma razón social aparece en los tres registros. \textbf{No prueba} que
el árido de la cantera «Macarena» sea el que se usa en esas obras: una empresa puede comprar
material a terceros, y el bacheo asfáltico usa insumos que no salen de un río.

\textbf{Tampoco hay nada irregular en ninguno de los tres roles.} La concesión se tramita ante
el Juzgado de Minas, la obra se adjudica por un procedimiento de contratación, y la
integración vertical entre extracción, transporte y obra es corriente en el rubro.

\textbf{Lo que el caso ilustra} es la forma del circuito, no una anomalía: \textbf{el material
del cauce, la empresa que lo mueve y la calle que se pavimenta con él están en la misma cadena,
y los dos extremos de esa cadena están en jurisdicciones distintas}. El departamento que aporta
el árido no recauda por la obra que se hace con él, y lo que pueda recaudar por canon depende de un convenio que este libro no obtuvo.

\textbf{Y una magnitud que conviene tener presente:} el capítulo \ref{cap:trabajo} establece que
la explotación de minas y canteras ocupa a \textbf{cincuenta y una personas} en todo el
departamento.
\end{cautela}

\section{Qué refutaría esta lectura}

Como toda inferencia de este libro, ésta tiene condiciones de falsación, y son concretas.

\textbf{Si el informe de dominio muestra que los propietarios de las grandes fincas son
caldereños}, y no gente de la capital ni sociedades con domicilio en ella, se cae buena parte
del argumento. La única explotación agropecuaria que concentra cinco mil ochocientas setenta
hectáreas —el cincuenta y uno por ciento de la superficie agropecuaria del departamento—
tiene un titular, y este libro no sabe quién es.

\textbf{Si el padrón de la empresa de agua muestra que la red efectivamente llega} y que el
problema es de conexión domiciliaria y no de tendido, entonces la carencia de servicios no
describe un territorio desatendido sino uno mal medido por este trabajo.

\textbf{Si las respuestas al pedido de derecho a réplica muestran que el municipio pidió y la
Provincia negó}, con expedientes y fechas, entonces el problema no es de diseño institucional
sino de decisiones políticas concretas, y tiene responsables identificables. Sería un libro
distinto, y probablemente uno mejor.

\textbf{Si aparecieran comunidades con ocupación anterior al archivo}, la lectura necesitaría un piso que no tiene. Esta condición \textbf{se cumplió en parte mientras este libro se escribía}: no aparecieron títulos comunitarios ni posesiones veinteñales, pero sí un reconocimiento estatal de ocupación. El capítulo \ref{cap:resistencias} documenta que hay dos comunidades kollas registradas
en el departamento ---Kondorwaira, con personería desde 2009 en Potrero de Castilla, y
Urkuwasi, en trámite--- y que el Equipo Técnico de la Ley 26.160 certificó en 2023 que
\textbf{la zona alta de la cuenca es territorio de ocupación actual, tradicional y pública} de
la primera, según relevamiento del INAI de 2011. En Potrero de Castilla, además, en
\textbf{septiembre de 1937} el Decreto 1421 negó la aquiescencia para instalar una escuela
nacional de la Ley 4874 alegando que el Consejo General de Educación «\textit{instalará
próximamente escuelas fiscales}» allí, y \textbf{veinticuatro días después} el Decreto 1500
derogó ese artículo y la concedió, porque la Inspección Nacional verificó el porcentaje de
analfabetos y comprobó que \textbf{no existía escuela provincial en la localidad ni en lugares
próximos}.

Eso no derrumba la lectura de este capítulo, pero le agrega un piso que no tenía.
\textbf{Debajo del dispositivo que el mercado encontró armado hay una ocupación anterior a
todo el archivo que este libro recorrió}. El cruce catastral del capítulo \ref{cap:resistencias} muestra que ningún loteo documentado toca todavía la parcela que la comunidad reclama; la pregunta que queda es qué ocurrirá en las fincas que la rodean.

Tres de las cuatro se verifican con documentos identificados por número en el apéndice \ref{cap:pedidos}; la tercera, con las respuestas al derecho a réplica. Ninguna
requiere una investigación nueva: requieren ir a buscarlas.

\section{Una última cosa}

En enero de 1910, los vecinos del pueblo de La Caldera le pidieron al gobierno de la
provincia que los defendiera «de las invasiones del Río del mismo nombre». El gobierno
asignó dos mil pesos, nombró una comisión de siete vecinos y le pidió al municipio que
aportara lo que pudiera. La obra se hizo y costó un cuarenta y ocho por ciento más de lo
previsto.

En febrero de 2026 el municipio seguía construyendo bateas y muros de contención.

Ciento dieciséis años, el mismo río, la misma respuesta: contener el agua. En todo ese
tiempo, ningún acto de los que este libro pudo leer se pregunta la otra cosa —qué se
construyó, dónde y con qué autorización sobre las márgenes que había que defender.

Esa pregunta la formula por primera vez una sentencia de 2021. Y es, para este libro, el
hecho institucional más importante de la historia del departamento: no porque haya resuelto
nada —el plan tardó cuatro años y el río volvió a desbordarse en el medio— sino porque
invirtió el orden de la pregunta. Durante un siglo se preguntó cómo defender al pueblo del
río. Recién ahora se pregunta qué se hizo con el río.

Mientras se cerraba este relevamiento ---resolución del 3 de septiembre de 2026, publicada el 10---, el Gobierno provincial dio de baja un contrato de doscientos dieciséis millones de pesos para construir defensas sobre el río La Caldera, \textbf{porque el cauce se había corrido y la solución proyectada ya no servía}. En 1940, un decreto había dejado constancia de que otras obras de defensa se suspendieron por la misma razón. \textbf{Ochenta y seis años, y el río sigue llegando antes que el expediente.}

Y hay un detalle final que resume el libro entero mejor que cualquier conclusión. Cuando la resolución que aprobó el Plan ordenó, en 2026, que el Plan de Manejo quedara a disposición de la comunidad, lo hizo en el domicilio que el colectivo actor había constituido (capítulo \ref{cap:amparo}): una casa de la avenida principal de La Caldera \textbf{identificada por el número de su medidor eléctrico}, porque la calle no tiene numeración. \textbf{El documento que ordena cómo se maneja la cuenca que abastece a la capital se consulta en una casa sin número.}

\begin{figure}[p]
\centering
\includegraphics[width=\textwidth,height=0.88\textheight,keepaspectratio]{fig-agua-actos.png}
\caption[El agua del departamento de La Caldera, 1885--1946]{\textbf{El agua del departamento
de La Caldera, 1885--1946.} \textbf{Arriba}, las fincas que nombra cada uno de los actos de
agua que el archivo registra, coloreadas por el año del acto, sobre los cursos del IGN y el
límite departamental; el círculo marca \textbf{las juntas del Wierna con el río La Caldera},
que es el punto desde el cual la Ley 11.078 de 1929 autorizó a la Municipalidad de Campo Santo
---hoy General Güemes--- a vigilar y prohibir bocas-tomas. \textbf{Abajo}, el arco completo en
seis actos: conceder (1885), reglamentar (1914), perder (1929), compartir (1931), inscribir
(1942), colaborar (1946). \textbf{Lo que no puede dibujarse}: el recorrido de la acequia de
Urquiza, porque el croquis que su escrito dice acompañar no se publicó, y las superficies
regadas, porque ningún acto las georreferencia; las fincas se ubican por el campo «finca» del
catastro vigente y no por el título. Elaboración propia; dibujo de
\texttt{scripts/mapa\_agua.py} del repositorio \texttt{mapas\_caldera} (apéndice
\ref{cap:fuentes}).}
\label{fig:agua}
\end{figure}
